\documentclass[11pt, a4paper]{article}

\usepackage[a4paper, margin=2.5cm]{geometry}
\usepackage{fontspec}
\usepackage[english, bidi=basic, provide=*]{babel}
\babelprovide[import, onchar=ids fonts]{english}
\babelfont{rm}{Noto Sans}

\usepackage{amsmath, amssymb, amsthm}
\usepackage{xcolor}
\usepackage{enumitem}
\usepackage{parskip}
\usepackage{hyperref}

\definecolor{accent}{HTML}{1E3A5F}
\definecolor{muted}{HTML}{4B5563}

\hypersetup{
    colorlinks=true,
    linkcolor=accent,
    urlcolor=accent,
    citecolor=accent
}

\newtheorem{theorem}{Theorem}
\renewcommand{\qedsymbol}{$\blacksquare$}

\setlist[itemize]{topsep=4pt, itemsep=3pt}
\setlist[enumerate]{topsep=4pt, itemsep=4pt}

\begin{document}

% Title Block
\begin{center}
    {\Large\bfseries\color{accent} Assignment: NP-Completeness of Zero-Swap Qubit Allocation (ZSQ-Alloc)\par}
    \vspace{0.4em}
    {\normalsize\color{muted} Course: Advances in Algorithms\par}
\end{center}
\vspace{0.2em}
\noindent{\color{accent}\rule{\linewidth}{0.8pt}}
\vspace{0.6em}

\section{Main Chosen Topic and Description of the Problem}

\noindent\textbf{Chosen Topic:} Hardware Compilation Strategy in Quantum Computing --- Zero-Swap Qubit Allocation (ZSQ-Alloc)

\subsection*{Context and Motivation}
Modern Noisy Intermediate-Scale Quantum (NISQ) processors (such as superconducting architectures) have limited physical qubit connectivity: a two-qubit entangling gate (e.g., CNOT) can only be executed between two physical qubits that share a direct hardware coupler. When compiling a logical quantum circuit onto a physical processor, the compiler must decide an initial mapping strategy that assigns logical (program) qubits to physical qubits. If a two-qubit gate acts on logical qubits mapped to non-adjacent physical qubits, the compiler must insert costly SWAP gates (each requiring 3 CNOT gates), which increases circuit depth and error rates.

A fundamental decision subproblem in quantum compilation strategy is determining whether a given quantum circuit can be mapped onto a target quantum architecture without inserting any SWAP gates.

\subsection*{Formal Problem Definition}
We define the Zero-Swap Qubit Allocation (ZSQ-Alloc) problem as a formal decision language $L_{\text{ZSQ}}$.

\begin{itemize}
    \item \textbf{Input Instance:} A tuple $\mathcal{I} = (P, \Psi, G_Q)$, where:
    \begin{itemize}
        \item $P = \{p_1, p_2, \dots, p_k\}$ is a finite set of logical (program) qubits.
        \item $\Psi = \{(p_{a_1}, p_{b_1}), \dots, (p_{a_m}, p_{b_m})\} \subseteq P \times P$ is the set of two-qubit entangling gate interactions (e.g., undirected CNOT pairs) required by the quantum circuit, with $p_{a_i} \neq p_{b_i}$.
        \item $G_Q = (Q, E_Q)$ is an undirected hardware coupling graph, where $Q = \{q_1, q_2, \dots, q_n\}$ is the set of physical qubits ($k \le n$) and $(q_u, q_v) \in E_Q$ if and only if the hardware supports a direct two-qubit gate between $q_u$ and $q_v$.
    \end{itemize}
    \item \textbf{Decision Question:} Does there exist an injective allocation function $\phi: P \to Q$ such that every two-qubit interaction in $\Psi$ is mapped to an adjacent pair of physical qubits in $G_Q$? Formally:
    \[
    L_{\text{ZSQ}} = \left\{ \langle P, \Psi, G_Q \rangle \;\middle|\;
    \begin{aligned}
        &\exists \text{ injective } \phi: P \to Q \text{ such that } \\
        &\forall (p_a, p_b) \in \Psi,\; (\phi(p_a), \phi(p_b)) \in E_Q
    \end{aligned}
    \right\}
    \]
\end{itemize}

\section{References}

\noindent\hangindent=1.5em
Garey, M. R., \& Johnson, D. S. (1979). \textit{Computers and intractability: A guide to the theory of NP-completeness}. W. H. Freeman and Company.\par

\noindent\hangindent=1.5em
Siraichi, M. Y., Santos, V. F. d., Collange, S., \& Pereira, F. M. Q. (2018). Qubit allocation. \textit{Proceedings of the 2018 International Symposium on Code Generation and Optimization (CGO 2018)}, 113--125. \url{https://doi.org/10.1145/3168822}\par

\noindent\hangindent=1.5em
Siraichi, M. Y., Santos, V. F. d., Collange, S., \& Pereira, F. M. Q. (2019). Qubit allocation as a combination of subgraph isomorphism and token swapping. \textit{Proceedings of the ACM on Programming Languages}, 3(OOPSLA), 120:1--120:29. \url{https://doi.org/10.1145/3360546}\par

\section{Proof that \texorpdfstring{$L_{\text{ZSQ}}$}{LZSQ} is NP-Complete}

\begin{theorem}
The language $L_{\text{ZSQ}}$ is NP-Complete.
\end{theorem}

To prove Theorem~1, we show that (3.1) $L_{\text{ZSQ}} \in \text{NP}$, and (3.2) a known NP-Complete problem, Subgraph Isomorphism ($L_{\text{SIP}}$), is polynomial-time many-one reducible to $L_{\text{ZSQ}}$ ($L_{\text{SIP}} \le_p L_{\text{ZSQ}}$).

\subsection{Show that \texorpdfstring{$L_{\text{ZSQ}} \in \text{NP}$}{LZSQ in NP}}

\begin{proof}
We construct a deterministic polynomial-time verifier $V(\mathcal{I}, \phi)$ that takes an instance $\mathcal{I} = (P, \Psi, G_Q)$ and a candidate certificate $\phi: P \to Q$ (represented as an array of $k = |P|$ physical qubit indices):
\begin{enumerate}
    \item \textbf{Injectivity Check:} Verify that $\phi(p_i) \in Q$ for all $p_i \in P$, and that $\phi(p_i) \neq \phi(p_j)$ for all $i \neq j$. Using a boolean visited array of size $n = |Q|$, this takes $O(k + n)$ time.
    \item \textbf{Adjacency Check:} For each two-qubit gate $(p_a, p_b) \in \Psi$, check whether the edge $(\phi(p_a), \phi(p_b))$ exists in $E_Q$. Using an adjacency matrix representation of $G_Q$ (constructed in $O(n^2)$ time), each gate lookup takes $O(1)$ time, for a total of $O(m)$ time where $m = |\Psi|$.
    \item \textbf{Decision:} Accept if both checks pass; otherwise, reject.
\end{enumerate}
The total verification time is $O(n^2 + k + m)$, which is polynomial in the input size $|\mathcal{I}|$. Since a valid certificate $\phi$ exists if and only if $\mathcal{I} \in L_{\text{ZSQ}}$, we conclude that $L_{\text{ZSQ}} \in \text{NP}$.
\end{proof}

\subsection{Show that a Known NP-Complete Problem is Reducible to \texorpdfstring{$L_{\text{ZSQ}}$}{LZSQ}}

We reduce from the canonical NP-Complete problem Subgraph Isomorphism ($L_{\text{SIP}}$) (Garey \& Johnson, 1979):
\[
L_{\text{SIP}} = \left\{ \langle G_1, G_2 \rangle \;\middle|\;
\begin{aligned}
    &G_1 = (V_1, E_1),\; G_2 = (V_2, E_2),\; \exists \text{ injective } f: V_1 \to V_2 \\
    &\text{such that } \forall (u, v) \in E_1,\; (f(u), f(v)) \in E_2
\end{aligned}
\right\}
\]

\subsubsection*{Polynomial-Time Reduction Function ($R$)}
Given an arbitrary instance $\langle G_1, G_2 \rangle$ of Subgraph Isomorphism where $G_1 = (V_1, E_1)$ and $G_2 = (V_2, E_2)$, we construct an instance $R(\langle G_1, G_2 \rangle) = \langle P, \Psi, G_Q \rangle$ of ZSQ-Alloc as follows:
\begin{enumerate}
    \item \textbf{Logical Qubits:} Create one logical qubit $p_u \in P$ for each vertex $u \in V_1$ (so $|P| = |V_1|$).
    \item \textbf{Quantum Circuit Interactions:} For each undirected edge $(u, v) \in E_1$, add a two-qubit interaction $(p_u, p_v)$ to $\Psi$ (so $|\Psi| = |E_1|$).
    \item \textbf{Hardware Coupling Graph:} Set the physical coupling graph directly to $G_Q = G_2$, where $Q = V_2$ and $E_Q = E_2$.
\end{enumerate}

\noindent\textbf{Time Complexity of Reduction:} Copying the vertex and edge sets of $G_1$ and $G_2$ takes $O(|V_1| + |E_1| + |V_2| + |E_2|)$ time, which is linear (and therefore polynomial) in the size of the Subgraph Isomorphism instance.

\subsubsection*{Proof of Correctness (Equivalence)}
\begin{proof}
We prove that $\langle G_1, G_2 \rangle \in L_{\text{SIP}} \iff \langle P, \Psi, G_Q \rangle \in L_{\text{ZSQ}}$.

\begin{itemize}
    \item \textbf{Forward Direction ($\Rightarrow$):}
    Assume $\langle G_1, G_2 \rangle \in L_{\text{SIP}}$. Then there exists an injective function $f: V_1 \to V_2$ such that for every $(u, v) \in E_1$, $(f(u), f(v)) \in E_2$. Define the qubit allocation $\phi: P \to Q$ by $\phi(p_u) = f(u)$ for every $p_u \in P$.
    \begin{itemize}
        \item \textit{Injectivity:} Since $f$ is injective, $\phi$ is injective.
        \item \textit{Hardware Compliance:} For every two-qubit gate $(p_u, p_v) \in \Psi$, by construction $(u, v) \in E_1$. Because $f$ preserves edges, $(\phi(p_u), \phi(p_v)) = (f(u), f(v)) \in E_2 = E_Q$.
    \end{itemize}
    Thus, $\phi$ is a valid zero-swap allocation, and $\langle P, \Psi, G_Q \rangle \in L_{\text{ZSQ}}$.

    \item \textbf{Reverse Direction ($\Leftarrow$):}
    Assume $\langle P, \Psi, G_Q \rangle \in L_{\text{ZSQ}}$. Then there exists an injective allocation $\phi: P \to Q$ such that for every $(p_u, p_v) \in \Psi$, $(\phi(p_u), \phi(p_v)) \in E_Q$. Define $f: V_1 \to V_2$ by $f(u) = \phi(p_u)$ for every $u \in V_1$.
    \begin{itemize}
        \item \textit{Injectivity:} Since $\phi$ is injective, $f$ is injective.
        \item \textit{Edge Preservation:} For every edge $(u, v) \in E_1$, by construction $(p_u, p_v) \in \Psi$. Because $\phi$ is a valid allocation, $(f(u), f(v)) = (\phi(p_u), \phi(p_v)) \in E_Q = E_2$.
    \end{itemize}
    Thus, $G_1$ is isomorphic to a subgraph of $G_2$, so $\langle G_1, G_2 \rangle \in L_{\text{SIP}}$.
\end{itemize}

Since $L_{\text{SIP}} \le_p L_{\text{ZSQ}}$, $L_{\text{ZSQ}}$ is NP-Hard. Combining this with Section~3.1 ($L_{\text{ZSQ}} \in \text{NP}$), we conclude that $L_{\text{ZSQ}}$ is NP-Complete.
\end{proof}

\end{document}
